\documentclass[uplatex,dvipdfmx]{jlreq}
\usepackage[fleqn,tbtags]{mathtools}
\usepackage{jlreq-deluxe}
\usepackage{geometry}
\usepackage{booktabs}
\usepackage{hyperref}
\usepackage[noalphabet]{pxchfon} % must be after otf package
\usepackage{stix2} %欧文＆数式フォント
\usepackage[fleqn,tbtags]{mathtools} % 数式関連 (w/ amsmath)
\usepackage{url}
\usepackage{here}
\geometry{left=20mm,right=20mm,top=25mm,bottom=25mm}

\title{作業ログ\\ \large プロジェクト名: [加速度センサーを用いた身体動作による輪唱制御システムの構築と視覚フィードバック効果]}
\author{報告者: 早川　和輝}
\date{報告日: 2026年5月27日}

\begin{document}
\maketitle

\section{進捗サマリー(要約)}
\begin{itemize}
    \item 全体進捗率: [15\%]
    \item マイルストーン達成状況: [達成]
    \item 主要成果:
    \begin{itemize}
        \item 成果1: 「かえるの合唱」の楽譜作成，および各音符の周波数(Hz)数値出力の実装．
        \item 成果2: Minimライブラリを利用したアコースティックバイオリンの音色合成プログラムの作成．
    \end{itemize}
    \item 重大課題(影響度:無し): なし
\end{itemize}

\section{作業ログ(詳細トラッキング)}
\begin{table}[H]
    \centering
    \small
    \begin{tabular}{lp{1.5cm}p{1.5cm}llp{2.5cm}p{3cm}}
        \toprule
        作業項目 & 開始日 & 完了日 & 状態 & 進捗率 & 依存関係・ブロッカー & コメント \\
        \midrule
        楽譜作成・Hz出力 & [5/23] & [5/23] & 完了 & 100\% & なし & 1オクターブ高い周波数配列(melody)を作成完了． \\
        バイオリンの音色作成 & [5/25] & [5/27] & 完了 & 100\% & なし &  AcousticViolinInstrumentの実装．波形とフィルターを調整． \\
        \bottomrule
    \end{tabular}
\end{table}

\section{担当箇所の計画前のGCと現在のGC}
担当箇所である「楽器演奏実装」における，当初計画時（計画前）のガントチャート（GC）の予定と，現在の進捗・変更状況の比較である．
\begin{table}[H]
    \centering
    \small
    \begin{tabular}{lp{4cm}p{4cm}l}
        \toprule
        タスク名 & 計画前のGC（予定期間） & 現在のGC（実績・修正期間） & 進捗状態 \\
        \midrule
        楽譜作成・Hz出力 & [5/23 $\sim$ 5/25] & [5/23 $\sim$ 5/23] & 完了 \\
        バイオリンの音色作成 & [5/29 $\sim$ 6/6] & [5/25 $\sim$ 5/27] & 完了 \\
        ギロの音色作成 & [5/29 $\sim$ 6/6] & [5/27 $\sim$ 6/6] & 未着手 \\
        ArudinoとProcessingの連携 & [6/3 $\sim$ 6/7] & [6/3 $\sim$ 6/7] & 未着手 \\
        音の動作確認 & [6/5 $\sim$ 6/9] & [6/5 $\sim$ 6/9] & 未着手 \\
        テンポを変更させる処理の追加 & [6/8 $\sim$ 6/13] & [6/8 $\sim$ 6/13] & 未着手 \\
        同期との連携 & [6/11 $\sim$ 6/17] & [6/11 $\sim$ 6/17] & 未着手 \\
        \bottomrule
    \end{tabular}
\end{table}
\section{日ごとの作業時間と作業内容}
今週実施した日ごとの具体的な作業時間および対応内容の詳細です．
\begin{table}[H]
    \centering
    \small
    \begin{tabular}{lll}
        \toprule
        日付 & 作業時間 & 作業内容 \\
        \midrule
        {}[5/23] & [0.5時間] &「かえるの合唱」の楽譜データの作成，および半拍ごとのHz数値の割り出し \\
        {}[5/25] & [1.0時間] & Processing環境およびMinimライブラリのセットアップ，発音の基本検証 \\
        {}[5/25] & [1.5時間] & バイオリン音色クラス（AcousticViolinInstrument）の実装  \\
        {}[5/26] & [1.0時間] & ノイズ・フィルター（MoogFilter）の調整，およびFFTによるスペクトラム描画の実装  \\
        {}[5/27] & [1.0時間] & 作成したバイオリン音色の全体の動作チェックと音質調整，および本作業ログの作成  \\
        \midrule
        \textbf{合計} & \textbf{[5.0時間]} & \\
        \bottomrule
    \end{tabular}
\end{table}
\section{課題管理}
\begin{table}[H]
    \centering
    \small
    \begin{tabular}{p{3cm}p{3cm}llp{3cm}ll}
        \toprule
        課題 & 発生原因 & 影響度 & 優先度 & 対策 & 期限 & 状態 \\
        \midrule
        擦弦音のリアルさ向上 & 単純な波形のみでは機械音になる & 高 & 高 & ピンクノイズとMoogFilterを追加し摩擦音を表現 & [5/27] & 解決 \\
        バイオリンの音色の特徴の尊重 & 低い音程だとバイオリンの良さが損なわれる & 高 & 高 & １オクターブ高い周波数に変換 & [5/27] & 解決 \\
        \bottomrule
    \end{tabular}
\end{table}
\vspace{0.5em}
\noindent ※影響度=プロジェクトへのインパクト，優先度=対応の緊急性

\section{AI 利用状況と検証(再現性重視)}
\subsection{利用内容}
\begin{itemize}
    \item 利用目的: \LaTeX による作業ログのテンプレート作成
    \item  入力(プロンプト要旨):
    \begin{itemize}
        \item 指定された作業ログのテンプレート項目の生成を依頼 ．
    \end{itemize}
    \item  出力概要:
    \begin{itemize}
        \item テンプレートの構成のコンパイル可能な作業ログの\LaTeX ソースコード ．
    \end{itemize}
\end{itemize}

\subsection{検証方法}
\begin{itemize}
    \item  検証手法: 生成された\LaTeX ソースコードの構造確認およびコンパイル検証
    \item  参照資料: 提供された作業ログテンプレートPDF
    \item  検証手順:
    \begin{enumerate}
        \item 生成されたコードのセクション構成が，テンプレート項目を漏れなく満たしているか目視で確認．
        \item \LaTeX 環境において実際にコンパイルを行い，エラーを吐かずに正しくPDFが生成されるかを確認．
    \end{enumerate}
\end{itemize}

\section{次回計画}
\begin{itemize}
    \item  目標: [バイオリンの音色のさらなる改善と，ギロの音色作成と]
    \item  達成基準: [バイオリンの音色がよりリアルになり，ギロの音色が完成すること]
    \item  期限: [6/6]
\end{itemize}

\section{その他}
連絡事項や，他メンバーとの連携が必要な懸念点があれば記載


\end{document}